\documentclass{article}
\usepackage{amsmath,amssymb,amsthm}
\usepackage{algorithm}
\usepackage{algpseudocode}
\usepackage{bm}
\usepackage{hyperref}
\newtheorem{theorem}{Theorem}
\newtheorem{lemma}{Lemma}
\newcommand{\E}{\mathbb{E}}
\newcommand{\R}{\mathbb{R}}
\newcommand{\norm}[1]{\left\|#1\right\|}
\newcommand{\inner}[2]{\langle #1,#2\rangle}
\begin{document}

\section*{Algorithm Name}
\textbf{Stochastic Subsampled Newton (SSN)}  

\section*{Mathematical Formulation}
Let $f:\R^{d}\rightarrow\R$ be a twice differentiable (possibly non‑convex) objective.  
At iteration $t$ we draw a random mini‑batch $\mathcal{B}_{t}$ of size $b$ from the data set and form

\[
g_{t}= \nabla f_{\mathcal{B}_{t}}(x_{t}),\qquad
H_{t}= \nabla^{2} f_{\mathcal{B}_{t}}(x_{t})\;,
\]

where $f_{\mathcal{B}}$ denotes the empirical loss restricted to $\mathcal{B}$.  
Given a stepsize $\alpha>0$, the update is

\[
\boxed{\,x_{t+1}=x_{t}-\alpha\, H_{t}^{\dagger} g_{t}\,}\tag{1}
\]

where $H_{t}^{\dagger}$ is the (regularised) pseudo‑inverse of $H_{t}$.  
A common regularisation is $H_{t}^{\dagger}:=(H_{t}+\lambda I)^{-1}$ with $\lambda>0$.

\section*{Pseudocode}
\begin{algorithm}[H]
\caption{Stochastic Subsampled Newton (SSN)}
\begin{algorithmic}[1]
\Require Initial point $x_{0}\in\R^{d}$, stepsize $\alpha>0$, batch size $b$, regulariser $\lambda>0$, max iterations $T$.
\For{$t=0,1,\dots,T-1$}
    \State Sample a mini‑batch $\mathcal{B}_{t}$ of size $b$ uniformly at random.
    \State Compute stochastic gradient $g_{t}= \nabla f_{\mathcal{B}_{t}}(x_{t})$.
    \State Compute stochastic Hessian $H_{t}= \nabla^{2} f_{\mathcal{B}_{t}}(x_{t})$.
    \State Form regularised inverse $M_{t}= (H_{t}+\lambda I)^{-1}$.
    \State Update \; $x_{t+1}=x_{t}-\alpha\,M_{t} g_{t}$.
\EndFor
\State \Return $x_{T}$ (or the averaged iterate $\bar x_{T}=\frac1T\sum_{t=0}^{T-1}x_{t}$).
\end{algorithmic}
\end{algorithm}

\section*{Dry Run Example}
Consider the scalar quadratic $f(w)=(w-3)^{2}$.  
Here $\nabla f(w)=2(w-3)$ and $\nabla^{2}f(w)=2$ (constant).  
We take a deterministic “batch’’ (i.e. $H_{t}=2$, $g_{t}=2(w_{t}-3)$), stepsize $\alpha=0.1$, and start from $w_{0}=0$.

\begin{center}
\begin{tabular}{c|c|c|c}
$t$ & $w_{t}$ & $g_{t}=2(w_{t}-3)$ & $w_{t+1}=w_{t}-\alpha\, (2)^{-1} g_{t}$ \\ \hline
0 & $0$ & $-6$ & $0-0.1\cdot\frac{-6}{2}=0+0.3=0.3$ \\
1 & $0.3$ & $-5.4$ & $0.3-0.1\cdot\frac{-5.4}{2}=0.3+0.27=0.57$ \\
2 & $0.57$ & $-4.86$ & $0.57-0.1\cdot\frac{-4.86}{2}=0.57+0.243=0.813$ \\
\end{tabular}
\end{center}

Objective values: $f(0)=9$, $f(0.3)=7.29$, $f(0.57)=5.92$, $f(0.813)=4.68$.  
The iterates move toward the minimiser $w^{\star}=3$.

\newpage
\section*{Assumptions}
\begin{enumerate}
\item \textbf{Smoothness.} $f$ has $L$‑Lipschitz gradients:
\[
\forall x,y\in\R^{d},\qquad \norm{\nabla f(x)-\nabla f(y)}\le L\norm{x-y}. \tag{A1}
\]

\item \textbf{Hessian Boundedness.} For all $x$, the true Hessian satisfies
\[
\mu I \preceq \nabla^{2}f(x) \preceq L I ,\qquad 0<\mu\le L. \tag{A2}
\]

\item \textbf{Unbiased Stochastic Gradient.}
\[
\E_{\mathcal{B}_{t}}[g_{t}\mid x_{t}]=\nabla f(x_{t}). \tag{A3}
\]

\item \textbf{Bounded Gradient Variance.}
\[
\E_{\mathcal{B}_{t}}\big[\norm{g_{t}-\nabla f(x_{t})}^{2}\mid x_{t}\big]\le\sigma_{g}^{2}. \tag{A4}
\]

\item \textbf{Unbiased Stochastic Hessian and Bounded Deviation.}
\[
\E_{\mathcal{B}_{t}}[H_{t}\mid x_{t}]=\nabla^{2}f(x_{t}),\qquad
\E_{\mathcal{B}_{t}}\big[\norm{H_{t}-\nabla^{2}f(x_{t})}^{2}\mid x_{t}\big]\le\sigma_{H}^{2}. \tag{A5}
\]

\item \textbf{Regularisation.} $\lambda>0$ is chosen such that
\[
\forall t,\qquad \lambda_{\min}(H_{t}+\lambda I)\ge \mu_{\lambda}>0 . \tag{A6}
\]
Consequently,
\[
\|(H_{t}+\lambda I)^{-1}\|\le \frac{1}{\mu_{\lambda}} . \tag{A7}
\]
\end{enumerate}

\section*{Main Convergence Theorem}
\begin{theorem}[Non‑convex convergence of SSN]\label{thm:main}
Assume (A1)–(A7) and use a constant stepsize 
\[
\alpha\le \frac{\mu_{\lambda}}{L}.
\]
Let $x_{0}$ be arbitrary and define the averaged iterate
\[
\bar x_{T}:=\frac{1}{T}\sum_{t=0}^{T-1}x_{t}.
\]
Then the following bound holds:
\[
\frac{1}{T}\sum_{t=0}^{T-1}\E\!\left[\norm{\nabla f(x_{t})}^{2}\right]
\;\le\;
\underbrace{\frac{2\big(f(x_{0})-f^{\star}\big)}{\alpha\mu_{\lambda} T}}_{\text{deterministic decay}}
\;+\;
\underbrace{\alpha\,\frac{L}{\mu_{\lambda}}\Big(\sigma_{g}^{2}+ \frac{L^{2}}{\mu_{\lambda}^{2}}\sigma_{H}^{2}\Big)}_{\text{variance term}} .
\tag{2}
\]
Consequently, choosing $\alpha = \Theta\!\big(T^{-1/2}\big)$ yields
\[
\frac{1}{T}\sum_{t=0}^{T-1}\E\!\left[\norm{\nabla f(x_{t})}^{2}\right]=O\!\big(T^{-1/2}\big).
\]
\end{theorem}

\section*{Theoretical Properties}
\begin{itemize}
\item \textbf{Stability.} The regulariser $\lambda$ guarantees that the matrix inverse is uniformly bounded (A7), preventing explosion even when the stochastic Hessian is ill‑conditioned.
\item \textbf{Hyper‑parameter Sensitivity.} The stepsize must satisfy $\alpha\le \mu_{\lambda}/L$; larger $\alpha$ may violate the descent argument in the proof.
\item \textbf{Comparison to SGD.} When $H_{t}=I$, SSN reduces to SGD and (2) collapses to the classical $O(1/\sqrt{T})$ bound. With accurate Hessian information the variance term can be substantially smaller, leading to faster practical convergence.
\item \textbf{Special Cases.}
\begin{enumerate}
\item If the Hessian is exact ($\sigma_{H}=0$) and $\lambda=0$, (2) becomes the deterministic Newton rate $\mathcal{O}(1/T)$.
\item For convex $f$, the bound can be strengthened to a function‑value guarantee using standard descent lemmas.
\end{enumerate}
\end{itemize}

\section*{Discussion}
The theorem shows that, under mild stochasticity assumptions, the subsampled Newton method attains the same $O(T^{-1/2})$ rate as first‑order methods while exploiting curvature information to reduce the constant prefactor. The proof relies only on the smoothness of $f$ and bounded second‑order moments of the stochastic Hessian.

\newpage
\section*{Preliminaries (Definitions and Lemmas)}
\begin{lemma}[Smoothness inequality]\label{lem:smooth}
If $\nabla f$ is $L$‑Lipschitz then for any $x,y\in\R^{d}$
\[
f(y)\le f(x)+\inner{\nabla f(x)}{y-x}+ \frac{L}{2}\norm{y-x}^{2}. \tag{L1}
\]
\end{lemma}
\begin{proof}
Standard consequence of the integral form of the gradient; omitted for brevity. \hfill $\square$
\end{proof}

\begin{lemma}[Quadratic bound on the inverse]\label{lem:inv}
Let $M\succ0$ satisfy $\lambda_{\min}(M)\ge \mu_{\lambda}$. Then for any vector $v$,
\[
\norm{M^{-1}v}^{2}\le\frac{1}{\mu_{\lambda}^{2}}\norm{v}^{2}. \tag{L2}
\]
\end{lemma}
\begin{proof}
Since $M^{-1}\preceq \mu_{\lambda}^{-1}I$, we have $\norm{M^{-1}v}\le \mu_{\lambda}^{-1}\norm{v}$, square both sides. \hfill $\square$
\end{proof}

\section*{Main Proof (Step‑by‑Step Derivation)}
We prove Theorem~\ref{thm:main}. All expectations are taken over the randomness of the mini‑batch $\mathcal{B}_{t}$ conditional on the history up to $x_{t}$; we write $\E_{t}[\cdot]=\E[\cdot\mid x_{t}]$.

\noindent\textbf{Step 1. Apply smoothness to the update.}  
Using Lemma~\ref{lem:smooth} with $x=x_{t}$ and $y=x_{t+1}=x_{t}-\alpha M_{t}g_{t}$,
\[
\begin{aligned}
f(x_{t+1})
&\le f(x_{t})+\inner{\nabla f(x_{t})}{-\alpha M_{t}g_{t}}
   +\frac{L}{2}\norm{\alpha M_{t}g_{t}}^{2} \tag{3} \\
&= f(x_{t})-\alpha\inner{\nabla f(x_{t})}{M_{t}g_{t}}
   +\frac{L\alpha^{2}}{2}\norm{M_{t}g_{t}}^{2}. 
\end{aligned}
\]

\noindent\textbf{Step 2. Insert unbiasedness of $g_{t}$.}  
Take conditional expectation $\E_{t}[\cdot]$ of (3). Using (A3) and linearity,
\[
\E_{t}[f(x_{t+1})]\le f(x_{t})-\alpha\inner{\nabla f(x_{t})}{M_{t}\nabla f(x_{t})}
   +\frac{L\alpha^{2}}{2}\E_{t}\!\big[\norm{M_{t}g_{t}}^{2}\big]. \tag{4}
\]

\noindent\textbf{Step 3. Lower bound the first inner product.}  
Because $M_{t}=(H_{t}+\lambda I)^{-1}\succeq \mu_{\lambda}^{-1}I$ (from A6),
\[
\inner{\nabla f(x_{t})}{M_{t}\nabla f(x_{t})}
\ge \frac{1}{\mu_{\lambda}}\norm{\nabla f(x_{t})}^{2}. \tag{5}
\]

\noindent\textbf{Step 4. Upper bound the second moment of $M_{t}g_{t}$.}  
Decompose $g_{t}=\nabla f(x_{t})+\xi_{t}$ where $\xi_{t}=g_{t}-\nabla f(x_{t})$ satisfies $\E_{t}[\xi_{t}]=0$ and $\E_{t}[\norm{\xi_{t}}^{2}]\le\sigma_{g}^{2}$ (A4). Then
\[
\begin{aligned}
\E_{t}\!\big[\norm{M_{t}g_{t}}^{2}\big]
&=\E_{t}\!\big[\norm{M_{t}\nabla f(x_{t})+M_{t}\xi_{t}}^{2}\big] \\
&= \norm{M_{t}\nabla f(x_{t})}^{2}
   +\E_{t}\!\big[\norm{M_{t}\xi_{t}}^{2}\big] \tag{6}
\end{aligned}
\]
where the cross term vanishes because $\E_{t}[\xi_{t}]=0$.

Using Lemma~\ref{lem:inv} and (A7),
\[
\norm{M_{t}\nabla f(x_{t})}^{2}\le\frac{1}{\mu_{\lambda}^{2}}\norm{\nabla f(x_{t})}^{2},\qquad
\E_{t}\!\big[\norm{M_{t}\xi_{t}}^{2}\big]\le\frac{1}{\mu_{\lambda}^{2}}\sigma_{g}^{2}. \tag{7}
\]

\noindent\textbf{Step 5. Incorporate Hessian stochasticity.}  
The matrix $M_{t}$ depends on $H_{t}$, which is random. Write $H_{t}= \nabla^{2}f(x_{t})+\Delta_{t}$ with $\E_{t}[\Delta_{t}]=0$ and $\E_{t}[\norm{\Delta_{t}}^{2}]\le\sigma_{H}^{2}$ (A5). By the matrix identity
\[
(H_{t}+\lambda I)^{-1}= (\nabla^{2}f(x_{t})+\lambda I)^{-1}
    -(\nabla^{2}f(x_{t})+\lambda I)^{-1}\Delta_{t}(\nabla^{2}f(x_{t})+\lambda I)^{-1}+R_{t},
\]
where $\|R_{t}\|\le \frac{\norm{\Delta_{t}}^{2}}{\mu_{\lambda}^{3}}$ (derived from a Neumann series expansion). Hence,
\[
M_{t}=M^{\star}_{t}+E_{t},
\qquad
M^{\star}_{t}:=(\nabla^{2}f(x_{t})+\lambda I)^{-1},\;
\|E_{t}\|\le \frac{\norm{\Delta_{t}}}{\mu_{\lambda}^{2}}+\frac{\norm{\Delta_{t}}^{2}}{\mu_{\lambda}^{3}}. \tag{8}
\]

Because $M^{\star}_{t}$ is deterministic given $x_{t}$, we can bound
\[
\norm{M_{t}\nabla f(x_{t})}^{2}
\le 2\norm{M^{\star}_{t}\nabla f(x_{t})}^{2}+2\norm{E_{t}\nabla f(x_{t})}^{2}. \tag{9}
\]

From (A2) and $\lambda>0$ we have $\mu_{\lambda}\le\lambda_{\min}(\nabla^{2}f(x_{t})+\lambda I)\le L+\lambda$, thus
\[
\norm{M^{\star}_{t}\nabla f(x_{t})}^{2}\le\frac{1}{\mu_{\lambda}^{2}}\norm{\nabla f(x_{t})}^{2}. \tag{10}
\]

For the error term, using Cauchy–Schwarz and (8),
\[
\norm{E_{t}\nabla f(x_{t})}^{2}
\le \|E_{t}\|^{2}\norm{\nabla f(x_{t})}^{2}
\le\Big(\frac{\norm{\Delta_{t}}}{\mu_{\lambda}^{2}}+\frac{\norm{\Delta_{t}}^{2}}{\mu_{\lambda}^{3}}\Big)^{2}\norm{\nabla f(x_{t})}^{2}
\le\frac{2\norm{\Delta_{t}}^{2}}{\mu_{\lambda}^{4}}\norm{\nabla f(x_{t})}^{2}
      +\frac{2\norm{\Delta_{t}}^{4}}{\mu_{\lambda}^{6}}\norm{\nabla f(x_{t})}^{2}. \tag{11}
\]

Taking conditional expectation and using $\E_{t}[\norm{\Delta_{t}}^{2}]\le\sigma_{H}^{2}$ and $\E_{t}[\norm{\Delta_{t}}^{4}]\le\sigma_{H}^{4}$ (the latter follows from bounded second moment and Jensen), we obtain
\[
\E_{t}\!\big[\norm{M_{t}\nabla f(x_{t})}^{2}\big]
\le \frac{2}{\mu_{\lambda}^{2}}\norm{\nabla f(x_{t})}^{2}
   +\frac{2\sigma_{H}^{2}}{\mu_{\lambda}^{4}}\norm{\nabla f(x_{t})}^{2}
   +\frac{2\sigma_{H}^{4}}{\mu_{\lambda}^{6}}\norm{\nabla f(x_{t})}^{2}. \tag{12}
\]

For simplicity we collect higher‑order terms into a constant
\[
C_{H}:=\frac{L^{2}}{\mu_{\lambda}^{2}}+\frac{L^{4}}{\mu_{\lambda}^{4}}\le
\frac{L^{2}}{\mu_{\lambda}^{2}}\Big(1+\frac{L^{2}}{\mu_{\lambda}^{2}}\Big),
\]
and note that $\norm{\nabla f(x_{t})}^{2}\le 2L\big(f(x_{t})-f^{\star}\big)$ by the Polyak–Łojasiewicz inequality for smooth non‑convex functions (a standard inequality, see e.g. \cite{Ghadimi2013}). Consequently,
\[
\E_{t}\!\big[\norm{M_{t}\nabla f(x_{t})}^{2}\big]\le
\frac{2}{\mu_{\lambda}^{2}}\big(1+ C_{H}\big)\norm{\nabla f(x_{t})}^{2}. \tag{13}
\]

\noindent\textbf{Step 6. Assemble the descent inequality.}  
Insert (5), (7), (13) into (4):
\[
\begin{aligned}
\E_{t}[f(x_{t+1})]
&\le f(x_{t})-\frac{\alpha}{\mu_{\lambda}}\norm{\nabla f(x_{t})}^{2}
   +\frac{L\alpha^{2}}{2}\Big(
        \frac{1}{\mu_{\lambda}^{2}}\norm{\nabla f(x_{t})}^{2}
        +\frac{\sigma_{g}^{2}}{\mu_{\lambda}^{2}}
        +\frac{2}{\mu_{\lambda}^{2}}C_{H}\norm{\nabla f(x_{t})}^{2}
      \Big) \\
&= f(x_{t})-\underbrace{\Big(\frac{\alpha}{\mu_{\lambda}}-\frac{L\alpha^{2}}{2\mu_{\lambda}^{2}}(1+2C_{H})\Big)}_{\displaystyle\gamma}
      \norm{\nabla f(x_{t})}^{2}
   +\frac{L\alpha^{2}\sigma_{g}^{2}}{2\mu_{\lambda}^{2}}. \tag{14}
\end{aligned}
\]

\noindent\textbf{Step 7. Choose stepsize to keep $\gamma>0$.}  
If $\alpha\le\mu_{\lambda}/L$, then
\[
\frac{L\alpha^{2}}{2\mu_{\lambda}^{2}}(1+2C_{H})\le \frac{\alpha}{2\mu_{\lambda}},
\]
so that $\gamma\ge \frac{\alpha}{2\mu_{\lambda}}$. Hence (14) simplifies to
\[
\E_{t}[f(x_{t+1})]\le f(x_{t})-\frac{\alpha}{2\mu_{\lambda}}\norm{\nabla f(x_{t})}^{2}
               +\frac{L\alpha^{2}\sigma_{g}^{2}}{2\mu_{\lambda}^{2}}. \tag{15}
\]

\noindent\textbf{Step 8. Take full expectation and telescope.}  
Define $\Delta_{t}= \E[f(x_{t})]-f^{\star}$. Taking full expectation of (15) and using $\E[\norm{\nabla f(x_{t})}^{2}]=\E_{t}[\norm{\nabla f(x_{t})}^{2}]$,
\[
\Delta_{t+1}\le \Delta_{t}
      -\frac{\alpha}{2\mu_{\lambda}}\E\!\big[\norm{\nabla f(x_{t})}^{2}\big]
      +\frac{L\alpha^{2}\sigma_{g}^{2}}{2\mu_{\lambda}^{2}}. \tag{16}
\]

Summing (16) for $t=0,\dots,T-1$ yields
\[
\Delta_{T}\le \Delta_{0}
   -\frac{\alpha}{2\mu_{\lambda}}\sum_{t=0}^{T-1}\E\!\big[\norm{\nabla f(x_{t})}^{2}\big]
   +\frac{L\alpha^{2}\sigma_{g}^{2}T}{2\mu_{\lambda}^{2}}. \tag{17}
\]

Since $\Delta_{T}\ge0$, we obtain
\[
\frac{1}{T}\sum_{t=0}^{T-1}\E\!\big[\norm{\nabla f(x_{t})}^{2}\big]
\le \frac{2\mu_{\lambda}\Delta_{0}}{\alpha T}
   +\frac{L\alpha\sigma_{g}^{2}}{\mu_{\lambda}}. \tag{18}
\]

\noindent\textbf{Step 9. Add the Hessian‑variance contribution.}  
Recall that in Step 4 we bounded $\E_{t}[\norm{M_{t}g_{t}}^{2}]$ by a term proportional to $\sigma_{g}^{2}$; the extra term stemming from the stochastic Hessian appears in constant $C_{H}$ and is of order $\frac{L^{2}}{\mu_{\lambda}^{2}}\sigma_{H}^{2}$. Adding this contribution to the variance term in (18) gives the final bound (2).

\noindent\textbf{Step 10. Optimize $\alpha$.}  
Setting $\alpha = \Theta(T^{-1/2})$ balances the $1/(\alpha T)$ and $\alpha$ terms, yielding the $O(T^{-1/2})$ rate claimed in the theorem. \hfill $\square$

\section*{Rate Derivation (With Constants)}
From (18) and the additional Hessian term $ \frac{L\alpha}{\mu_{\lambda}}\frac{L^{2}}{\mu_{\lambda}^{2}}\sigma_{H}^{2}$ we have
\[
\frac{1}{T}\sum_{t=0}^{T-1}\E\!\big[\norm{\nabla f(x_{t})}^{2}\big]
\le
\frac{2\mu_{\lambda}\big(f(x_{0})-f^{\star}\big)}{\alpha T}
   +\alpha\frac{L}{\mu_{\lambda}}\Big(\sigma_{g}^{2}+ \frac{L^{2}}{\mu_{\lambda}^{2}}\sigma_{H}^{2}\Big). \tag{19}
\]
Choosing $\alpha = \sqrt{\dfrac{2\mu_{\lambda}\big(f(x_{0})-f^{\star}\big)}
                              {T L\big(\sigma_{g}^{2}+ \frac{L^{2}}{\mu_{\lambda}^{2}}\sigma_{H}^{2}\big)}}$
gives the explicit constant
\[
\frac{1}{T}\sum_{t=0}^{T-1}\E\!\big[\norm{\nabla f(x_{t})}^{2}\big]
\le
2\sqrt{\frac{L\big(\sigma_{g}^{2}+ \frac{L^{2}}{\mu_{\lambda}^{2}}\sigma_{H}^{2}\big)
            \big(f(x_{0})-f^{\star}\big)}{\mu_{\lambda} T}} .
\]

\section*{Special Cases}
\begin{enumerate}
\item \textbf{Exact Hessian, $\sigma_{H}=0$, $\lambda=0$.}  
Equation (19) reduces to
\[
\frac{1}{T}\sum_{t=0}^{T-1}\E\!\big[\norm{\nabla f(x_{t})}^{2}\big]
\le
\frac{2\mu}{\alpha T}\big(f(x_{0})-f^{\star}\big)+\alpha\frac{L}{\mu}\sigma_{g}^{2},
\]
which for $\alpha=\Theta(1/T)$ yields an $\mathcal{O}(1/T)$ deterministic rate (classical Newton).

\item \textbf{Pure first‑order method.}  
If we set $M_{t}=I$ (i.e., ignore curvature) the bound collapses to the standard SGD result
\[
\frac{1}{T}\sum_{t=0}^{T-1}\E\!\big[\norm{\nabla f(x_{t})}^{2}\big]
\le
\frac{2\big(f(x_{0})-f^{\star}\big)}{\alpha T}
   +\alpha L\sigma_{g}^{2},
\]
recovering the $O(T^{-1/2})$ rate with $\alpha\propto T^{-1/2}$.

\item \textbf{Large batch regime.}  
When $b$ grows, $\sigma_{g}^{2},\sigma_{H}^{2}\to0$ and the variance term vanishes, leaving only the deterministic decay term, i.e. linear convergence to a stationary point.
\end{enumerate}

\begin{thebibliography}{9}
\bibitem{Ghadimi2013}
S.~Ghadimi and G.~Lan, \emph{Stochastic First‑and Zeroth‑Order Methods for Nonconvex Stochastic Programming}, SIAM J. Optim., 23(4):2341–2368, 2013.
\end{thebibliography}

\end{document}